\begin{longtable}[c]{|c|l|m{0.4\textwidth}|}
    \caption{Endpoints de autenticación y gestión de conexiones}
    \label{tab:endpointsSprint01} \\
    
     \hline
    \rowcolor{gray!50}
    \multicolumn{1}{|c|}{\textbf{Método HTTP}} & \multicolumn{1}{c|}{\textbf{Endpoint}} & \multicolumn{1}{c|}{\textbf{Descripción}} \\
    \hline
    \endfirsthead

    \hline
    \rowcolor{gray!50}
    \multicolumn{1}{|c|}{\textbf{Método HTTP}} & \multicolumn{1}{c|}{\textbf{Endpoint}} & \multicolumn{1}{c|}{\textbf{Descripción}} \\
    \hline
    \endhead

    \multicolumn{2}{r}{\textit{Continúa en la siguiente página...}} \\
    \endfoot
    
    \endlastfoot

    GET & \texttt{/connections} & Indica todas las conexiones del usuario. \\
    \hline
    GET & \texttt{/connections/conn-id} & Indica los detalles de una conexión específica. \\
    \hline
    POST & \texttt{/connections} & Agrega una nueva conexión. \\
    \hline
    PUT & \texttt{/connections/conn-id} & Actualiza una conexión existente. \\
    \hline
    DELETE & \texttt{/connections/conn-id} & Elimina una conexión \\
    \hline
    POST & \texttt{/connections/conn-id/test} & Permite realizar pruebas de conexión. \\
    \hline
    POST & \texttt{/auth/register} & Permite registrar un nuevo usuario. \\
    \hline
    POST & \texttt{/auth/login} & Permite iniciar sesión y solicita la generación un token JWT. \\
    \hline
    
\end{longtable}